\chapter{The Ornstein--Uhlenbeck channel and Tweedie's identity}
\label{ch:ou-channel-tweedie}

This chapter builds the forward process of the paper from nothing. It writes down the
stochastic differential equation that adds the noise, solves it, and derives the three
facts on which every later chapter rests: the noised law is a Gaussian convolution of
the clean law; its per-site variance is one at every noise level; and its score is a
rescaled posterior mean (Tweedie's identity). It ends with the observation protocol of
the paper, one noised view per sequence on a twelve-level schedule, written out level by
level. Throughout, $\va\in\Reals^{\nsites}$ is a clean sequence with law $\Pz$,
$\vx_t$ its noised version at level $t$ with law $\Pt$, and $\vz\sim\normal{0}{\Id}$ a
standard Gaussian vector independent of $\va$. In derivations we abbreviate
$\alpha=e^{-t}$, as the paper's appendix does; in displayed statements we keep
$e^{-t}$, as the paper's body does.

\section{Stochastic differential equations in one paragraph}
\label{sec:ou:sde}

A \term{Brownian motion} $(W_t)_{t\ge0}$ is a random continuous path with $W_0=0$ whose
increments over disjoint intervals are independent and Gaussian,
$W_t-W_s\sim\normal{0}{t-s}$ for $s<t$. Its formal time derivative
$\eta(t)=\dd W_t/\dd t$ is \term{white noise}. It is \term{standardised} in the sense
that $\E[\eta(t)\eta(t')]=\delta(t-t')$, which is the same as saying that
$\int_s^t\eta(u)\dd u=W_t-W_s$ has variance $t-s$. A \term{stochastic differential
equation}
\begin{equation}
  \dd x_t=f(x_t,t)\,\dd t+g(t)\,\dd W_t
  \label{eq:ou:sde-generic}
\end{equation}
is shorthand for the recursion $x_{t+h}-x_t\approx f(x_t,t)\,h+g(t)\,(W_{t+h}-W_t)$
taken to the limit $h\to0$; $f$ is the \term{drift} (the deterministic push) and $g$ the
\term{diffusion coefficient} (the size of the random kick). The limit defines the It\^o
integral, and only two of its properties are needed here. First, for a deterministic
integrand $\phi$ the random variable $\int_0^t\phi(u)\dd W_u$ is Gaussian with mean zero
and variance $\int_0^t\phi(u)^2\dd u$ (the \term{It\^o isometry}). Second, a linear
equation can be solved with an integrating factor exactly as an ordinary differential
equation can, because multiplying by a smooth deterministic function introduces no
It\^o correction term. In $\nsites$ dimensions $W_t$ has $\nsites$ independent
coordinates and \cref{eq:ou:sde-generic} is read coordinate by coordinate.

\section{The variance-preserving Ornstein--Uhlenbeck process}
\label{sec:ou:process}

\begin{definition}[The channel]
\label{def:ou:channel}
The \term{Ornstein--Uhlenbeck channel} of the paper is the coordinatewise process
\begin{equation}
  \dd\vx_t=-\vx_t\,\dd t+\sqrt2\,\dd\vect{W}_t,\qquad \vx_0=\va\sim\Pz,
  \label{eq:ou:sde}
\end{equation}
with $\vect{W}$ an $\nsites$-dimensional Brownian motion independent of $\va$. In white-noise
form, $\dd\vx_t/\dd t=-\vx_t+\sqrt2\,\gwn(t)$ with $\gwn$ standardised.
\end{definition}

The drift $-\vx_t$ pulls every coordinate towards zero at unit rate; the kick
$\sqrt2\,\dd\vect{W}_t$ is the same for every coordinate and independent across
coordinates. Nothing couples the sites: the channel acts on each $a_i$ separately, and
this is what keeps the posterior factor graph a tree in the paper's Section 2.

\begin{proposition}[Solution of the channel]
\label{prop:ou:solution}
For every $t\ge0$,
\begin{equation}
  \vx_t=e^{-t}\va+\sqrt{\Dt}\,\vz,\qquad \Dt=1-e^{-2t},\qquad \vz\sim\normal{0}{\Id},
  \label{eq:ou:solution}
\end{equation}
with $\vz$ independent of $\va$. Conditionally on $\va$, $\vx_t\sim\normal{e^{-t}\va}{\Dt\Id}$.
\end{proposition}

\begin{proof}
Take one coordinate and drop the index. Put $y_t=e^{t}x_t$. The factor $e^{t}$ is
deterministic and smooth, so the ordinary product rule applies:
\[
  \dd y_t=e^{t}x_t\,\dd t+e^{t}\,\dd x_t
        =e^{t}x_t\,\dd t+e^{t}\bigl(-x_t\,\dd t+\sqrt2\,\dd W_t\bigr)
        =\sqrt2\,e^{t}\,\dd W_t .
\]
Integrating from $0$ to $t$ gives $y_t=a+\sqrt2\int_0^t e^{u}\dd W_u$, hence
\[
  x_t=e^{-t}a+\sqrt2\int_0^t e^{-(t-u)}\dd W_u .
\]
The integrand $\sqrt2\,e^{-(t-u)}$ is deterministic, so by the It\^o isometry the
integral is Gaussian with mean zero and variance
\[
  \int_0^t 2\,e^{-2(t-u)}\dd u=2\cdot\frac{1-e^{-2t}}{2}=1-e^{-2t}=\Dt .
\]
It is independent of $a$ because $W$ is. Writing the integral as $\sqrt{\Dt}\,z$ with
$z\sim\normal01$ gives the scalar claim; the coordinates use independent Brownian
motions, so the vector claim follows with $\vz\sim\normal{0}{\Id}$.
\end{proof}

\begin{remark}[Why $\sqrt2$]
\label{rem:ou:sqrt2}
The variance $v_t=\Var(x_t)$ of one coordinate obeys $\dd v_t/\dd t=-2v_t+2$, whose
fixed point is $v=1$. With the kick $\dd W_t$ instead of $\sqrt2\,\dd W_t$ the fixed
point would be $\tfrac12$ and one would have $\Dt=(1-e^{-2t})/2$. The paper's
convention (\code{paper/notation.tex}, item 4) fixes the noise as standardised and
carries the $\sqrt2$ explicitly so that $\Dt=1-e^{-2t}$ throughout, in the forward and
in the reverse equation. The variance-preserving SDE of Song and co-authors,
$\dd\vx=-\tfrac12\beta(t)\vx\,\dd t+\sqrt{\beta(t)}\,\dd\vect{W}$, is the same process
when $\beta\equiv2$; for a general schedule $\beta$ the paper's $t$ is
$\tfrac12\int_0^{t'}\beta(u)\dd u$, and nothing below changes.
\end{remark}

\begin{remark}[The semigroup property]
\label{rem:ou:semigroup}
Noising for time $t$ and then for time $s$ is noising for time $t+s$:
if $\vx_t=e^{-t}\va+\sqrt{\Dt}\vz$ and $\vx_{t+s}=e^{-s}\vx_t+\sqrt{\Delta_s}\vz'$ with
$\vz'$ fresh, then $\vx_{t+s}=e^{-(t+s)}\va+\bigl(e^{-s}\sqrt{\Dt}\vz+\sqrt{\Delta_s}\vz'\bigr)$,
and the bracket is Gaussian with variance
$e^{-2s}(1-e^{-2t})+1-e^{-2s}=1-e^{-2(t+s)}=\Delta_{t+s}$ per coordinate. This is the
Markov property of the channel, used in \cref{ch:reverse-sde-generation}.
\end{remark}

\section{Variance preservation and the Gaussian limit}
\label{sec:ou:variance}

\begin{proposition}[Second moments of the noised law]
\label{prop:ou:variance}
Let $\Sig_0=\Cov(\va)$. Then $\E[\vx_t]=e^{-t}\E[\va]$ and
\begin{equation}
  \Cov(\vx_t)=e^{-2t}\Sig_0+\Dt\Id=:\Sig_t .
  \label{eq:ou:sigmat}
\end{equation}
In particular, if $\Var(a_i)=1$ for every site, then $\Var(x_{t,i})=e^{-2t}+\Dt=1$ at
every $t$: the channel \term{preserves the variance}.
\end{proposition}

\begin{proof}
From \cref{eq:ou:solution}, $\va$ and $\vz$ are independent, so covariances add:
$\Cov(\vx_t)=e^{-2t}\Cov(\va)+\Dt\Cov(\vz)$.
\end{proof}

The paper's chain (Section 2) has $a_1\sim\normal01$ and innovation variance
$\ivar=1-\corr^2$, which gives $\Var(a_i)=1$ and $\Cov(a_i,a_j)=\corr^{|i-j|}$ for
every innovation law; hence every family shares the same $\Sig_0$ and the same $\Sig_t$.

\begin{proposition}[The Gaussian limit]
\label{prop:ou:limit}
As $t\to\infty$, $\Pt$ converges in distribution to $\gamma=\normal{0}{\Id}$, whatever
$\Pz$ is, provided $\E|\va|^2<\infty$.
\end{proposition}

\begin{proof}
For a bounded continuous $\phi$, $\E\phi(\vx_t)=\E\phi(e^{-t}\va+\sqrt{\Dt}\vz)$. The
term $e^{-t}\va$ tends to zero in probability by Chebyshev's inequality,
$\Prob(e^{-t}|\va|>\epsilon)\le e^{-2t}\E|\va|^2/\epsilon^2$, and $\sqrt{\Dt}\vz\to\vz$
almost surely. Dominated convergence then gives $\E\phi(\vx_t)\to\E\phi(\vz)$.
\end{proof}

The limit is the stationary law of the process: if $\va\sim\gamma$ then
$\vx_t\sim\normal{0}{e^{-2t}\Id+\Dt\Id}=\gamma$ for every $t$. The rate of convergence is
the subject of \cref{ch:attenuation-theorems}: the Hermite moment of degree $m$ of
$\Pt$ is exactly $e^{-mt}$ times that of $\Pz$, so the memory of $\Pz$ fades order by
order rather than all at once.

\section{The noised law as a Gaussian convolution}
\label{sec:ou:convolution}

Write $\gamma_{\Delta}$ for the density of $\normal{0}{\Delta\Id}$ on $\Reals^{\nsites}$,
\begin{equation}
  \gamma_{\Delta}(\vect{w})=(2\pi\Delta)^{-\nsites/2}\exp\Bigl(-\frac{|\vect{w}|^2}{2\Delta}\Bigr).
  \label{eq:ou:gaussdensity}
\end{equation}

\begin{proposition}[Density of the noised law]
\label{prop:ou:convolution}
For every $t>0$ the law $\Pt$ has the density
\begin{equation}
  \Pt(\vx)=\int_{\Reals^{\nsites}}\Pz(\dd\va)\,\gamma_{\Dt}(\vx-e^{-t}\va)
         =\E_{\Pz}\bigl[\gamma_{\Dt}(\vx-e^{-t}\va)\bigr],
  \label{eq:ou:convolution}
\end{equation}
which is strictly positive and infinitely differentiable in $\vx$, whatever $\Pz$ is
(it need not have a density). Every derivative may be taken under the integral sign.
\end{proposition}

\begin{proof}
Conditionally on $\va$, $\vx_t$ has density $\gamma_{\Dt}(\vx-e^{-t}\va)$ by
\cref{prop:ou:solution}; integrating over $\va$ gives \cref{eq:ou:convolution}.
Positivity is clear. For differentiation under the integral, every $\vx$-derivative of
$\gamma_{\Dt}(\vx-e^{-t}\va)$ is a polynomial in $(\vx-e^{-t}\va)/\Dt$ times the same
Gaussian, and is therefore bounded uniformly in $\va$ on every bounded set of $\vx$ by a
constant, which is integrable against the probability measure $\Pz$; dominated
convergence applies.
\end{proof}

Two readings of \cref{eq:ou:convolution} are used. As an operator statement,
$\Pt$ is the convolution of the dilated law (the law of $e^{-t}\va$) with the Gaussian
$\gamma_{\Dt}$. As a factor-graph statement, the integrand factorises over sites,
\begin{equation}
  \gamma_{\Dt}(\vx-e^{-t}\va)=\prod_{i=1}^{\nsites}\like(x_i\mid a_i),\qquad
  \like(x_i\mid a_i)\propto\exp\Bigl[-\frac{(x_i-e^{-t}a_i)^2}{2\Dt}\Bigr],
  \label{eq:ou:unary}
\end{equation}
so conditioning on $\vx$ attaches one \term{unary likelihood factor} $\like$ to each
site and adds no edge among the latent variables. The code builds exactly this table:
\code{log\_likelihood\_matrix} in \code{src/markov\_score/noising.py} returns
$\log\like(x_i\mid\gridpt{k})$ on the quadrature grid, shifted per row by a constant
so that no row underflows; the shift cancels in every normalised message and changes
nothing.

\section{The signal-to-noise ratio}
\label{sec:ou:snr}

\begin{definition}[Signal-to-noise ratio]
\label{def:ou:snr}
The \term{signal-to-noise ratio} of the channel at level $t$ is the ratio of the
signal variance to the noise variance in one coordinate of \cref{eq:ou:solution} for
unit-variance data,
\begin{equation}
  \mathrm{SNR}_t=\frac{e^{-2t}}{\Dt}=\frac{e^{-2t}}{1-e^{-2t}}=\frac{1}{e^{2t}-1}.
  \label{eq:ou:snr}
\end{equation}
\end{definition}

For small $t$, $\Dt=2t+O(t^2)$ and $\mathrm{SNR}_t\approx1/(2t)$, which diverges: the
observation is almost clean. For large $t$, $\mathrm{SNR}_t\approx e^{-2t}$: the signal
is a small perturbation of pure noise, and $e^{-2t}$ is the natural small parameter of
the expansion in \cref{ch:attenuation-theorems}. The same quantity appears there under
the name $b=e^{-2t}/\Dt$, together with $s=e^{-t}\vx/\Dt$; the reader should keep in
mind that $b$ is the SNR. Inverting \cref{eq:ou:snr}, $t=\tfrac12\log(1+1/\mathrm{SNR}_t)$,
so the diffusion time is a logarithmic clock on the SNR.

\begin{table}[t]
\centering
\caption{The twelve-level schedule of the paper (\code{experiments/frozen\_config.py},
$t_k=0.05\cdot r^{k-1}$ with $r=60^{1/11}\approx1.4509$), with the analytic values of the
decay factor, the noise variance and the signal-to-noise ratio at each level. All
entries are computed from the definitions; none is measured.}
\label{tab:ou:levels}
\begin{tabular}{@{}rrrrr@{}}
\toprule
$k$ & $t_k$ & $e^{-t_k}$ & $\Delta_{t_k}$ & $\mathrm{SNR}_{t_k}$ \\
\midrule
1  & 0.0500 & 0.9512 & 0.0952 & 9.51 \\
2  & 0.0725 & 0.9300 & 0.1350 & 6.40 \\
3  & 0.1053 & 0.9001 & 0.1898 & 4.27 \\
4  & 0.1527 & 0.8584 & 0.2632 & 2.80 \\
5  & 0.2216 & 0.8012 & 0.3580 & 1.79 \\
6  & 0.3215 & 0.7250 & 0.4743 & 1.11 \\
7  & 0.4665 & 0.6272 & 0.6066 & 0.648 \\
8  & 0.6769 & 0.5082 & 0.7417 & 0.348 \\
9  & 0.9821 & 0.3745 & 0.8597 & 0.163 \\
10 & 1.4250 & 0.2405 & 0.9422 & 0.061 \\
11 & 2.0676 & 0.1265 & 0.9840 & 0.016 \\
12 & 3.0000 & 0.0498 & 0.9975 & 0.0025 \\
\bottomrule
\end{tabular}
\end{table}

\section{Tweedie's identity}
\label{sec:ou:tweedie}

The \term{score} of the noised law is the gradient of its log-density,
$\score(\vx,t)=\nabla_{\vx}\log\Pt(\vx)$. Tweedie's identity says that it is a
rescaled residual of the posterior mean of the clean sequence. We write
$\pmean{\va}=\E[\va\mid\vx_t=\vx]$ for that posterior mean, in the bracket notation of
the paper.

\begin{theorem}[Tweedie's identity]
\label{thm:ou:tweedie}
For every $t>0$ and every $\vx\in\Reals^{\nsites}$,
\begin{equation}
  \score(\vx,t)=\nabla_{\vx}\log\Pt(\vx)=-\frac{\vx-e^{-t}\pmean{\va}}{\Dt},
  \qquad \pmean{\va}=\E[\va\mid\vx_t=\vx].
  \label{eq:ou:tweedie}
\end{equation}
Equivalently, $\pmean{\va}=e^{t}\bigl(\vx+\Dt\,\score(\vx,t)\bigr)$. No assumption on
$\Pz$ is needed beyond $\E|\va|<\infty$.
\end{theorem}

\begin{proof}
Scalar case first. By \cref{prop:ou:convolution},
$\Pt(x)=\int\Pz(\dd a)\,\gamma_{\Dt}(x-\alpha a)$ with $\alpha=e^{-t}$, and the
derivative may be taken inside:
\[
  \frac{\dd}{\dd x}\gamma_{\Dt}(x-\alpha a)=-\frac{x-\alpha a}{\Dt}\,\gamma_{\Dt}(x-\alpha a).
\]
Hence
\[
  \Pt'(x)=-\frac1{\Dt}\int\Pz(\dd a)\,(x-\alpha a)\,\gamma_{\Dt}(x-\alpha a)
         =-\frac1{\Dt}\Bigl[x\,\Pt(x)-\alpha\int\Pz(\dd a)\,a\,\gamma_{\Dt}(x-\alpha a)\Bigr].
\]
The posterior of $a$ given $x_t=x$ is, by Bayes' rule, the probability measure
$\Pz(\dd a)\,\gamma_{\Dt}(x-\alpha a)/\Pt(x)$, so the last integral equals
$\Pt(x)\,\E[a\mid x_t=x]$. Dividing by $\Pt(x)>0$,
\[
  \frac{\Pt'(x)}{\Pt(x)}=-\frac{x-\alpha\,\E[a\mid x_t=x]}{\Dt}.
\]
The integrability of $|a|$ under $\Pz$ is what makes the last integral finite. In
$\nsites$ dimensions, $\nabla_{\vx}\gamma_{\Dt}(\vx-\alpha\va)=-\Dt^{-1}(\vx-\alpha\va)\gamma_{\Dt}(\vx-\alpha\va)$
and the same three lines go through coordinate by coordinate.
\end{proof}

\begin{remark}[The classical forms]
\label{rem:ou:robbins}
Robbins (1956) wrote the identity for an empirical Bayes problem in which an observation
$z\sim\normal{\mu}{\sigma^2}$ has an unknown prior on $\mu$: $\E[\mu\mid z]=z+\sigma^2\,(\log f)'(z)$
with $f$ the marginal density of $z$. Efron (2011) named it Tweedie's formula. The
channel is this setting with $\mu=e^{-t}\va$ and $\sigma^2=\Dt$:
$\E[e^{-t}\va\mid\vx]=\vx+\Dt\score(\vx,t)$, which is the second form in
\cref{thm:ou:tweedie}. The code uses it in both directions: \code{score\_from\_mean} in
\code{src/markov\_score/theory.py} turns a posterior mean into a score, and
\code{denoising\_readout} in \code{src/markov\_score/reverse.py} turns a score into a
posterior mean.
\end{remark}

\begin{proposition}[Second-order Tweedie identity]
\label{prop:ou:tweedie2}
For $t>0$ the posterior covariance is also a function of the noised law:
\begin{equation}
  \Cov(\va\mid\vx_t=\vx)=\frac{\Dt}{e^{-2t}}\Bigl(\Id+\Dt\,\nabla^2_{\vx}\log\Pt(\vx)\Bigr),
  \label{eq:ou:tweedie2}
\end{equation}
so that $\operatorname{div}\score(\vx,t)=\tr\nabla^2\log\Pt(\vx)=-\nsites/\Dt+(e^{-2t}/\Dt^2)\tr\Cov(\va\mid\vx_t=\vx)$.
\end{proposition}

\begin{proof}
Differentiate \cref{eq:ou:tweedie} once more. With $m(\vx)=\pmean{\va}$,
$\nabla^2\log\Pt=-\Id/\Dt+(\alpha/\Dt)\nabla_{\vx}m$. Writing
$m_i(\vx)=\int a_i\,\Pz(\dd\va)\gamma_{\Dt}(\vx-\alpha\va)/\Pt(\vx)$ and differentiating
the quotient with respect to $x_j$,
\[
  \partial_j m_i=-\frac1{\Dt}\Bigl(\E[a_i(x_j-\alpha a_j)\mid\vx]-m_i\,(x_j-\alpha m_j)\Bigr)
               =\frac{\alpha}{\Dt}\bigl(\E[a_ia_j\mid\vx]-m_im_j\bigr)
               =\frac{\alpha}{\Dt}\Cov(a_i,a_j\mid\vx),
\]
because the terms in $x_j$ cancel. Substituting gives
$\nabla^2\log\Pt=-\Id/\Dt+(\alpha^2/\Dt^2)\Cov(\va\mid\vx)$, which is
\cref{eq:ou:tweedie2} rearranged.
\end{proof}

\why{The divergence of the score is what implicit score matching
(\cref{ch:score-matching}) has to compute, and \cref{eq:ou:tweedie2} says what it is:
up to constants, the trace of the posterior covariance. It also shows that the score
is a gradient field whose Jacobian is symmetric, a property a network's output need not
have.}

\begin{corollary}[The Gaussian case]
\label{cor:ou:gaussian}
If $\Pz=\normal{0}{\Sig_0}$ then $\Pt=\normal{0}{\Sig_t}$ with $\Sig_t$ as in
\cref{eq:ou:sigmat}, the score is linear, $\score(\vx,t)=-\Sig_t^{-1}\vx$, and the
posterior mean is $\pmean{\va}=e^{-t}\Sig_0\Sig_t^{-1}\vx$.
\end{corollary}

\begin{proof}
A Gaussian pushed through \cref{eq:ou:solution} stays Gaussian, with the covariance of
\cref{prop:ou:variance}. Then $\log\Pt(\vx)=-\tfrac12\vx^{\!\top}\Sig_t^{-1}\vx+\text{const}$
and $\score=-\Sig_t^{-1}\vx$. Inverting \cref{eq:ou:tweedie},
\[
  \pmean{\va}=\frac{\vx+\Dt\score}{\alpha}
             =\frac{(\Id-\Dt\Sig_t^{-1})\vx}{\alpha}
             =\frac{(\Sig_t-\Dt\Id)\Sig_t^{-1}\vx}{\alpha}
             =\frac{\alpha^2\Sig_0\Sig_t^{-1}\vx}{\alpha}
             =\alpha\,\Sig_0\Sig_t^{-1}\vx .\qedhere
\]
\end{proof}

This is the \term{linear minimum-mean-square estimator} (LMMSE) of $\va$ from $\vx$,
and the paper's Proposition 1 (Section 2) shows that it is exactly what Gaussian
message closure computes for \emph{any} innovation law with the same $\Sig_0$. The
score $\Sgauss(\vx,t)=-\Sig_t^{-1}\vx$ of this \term{covariance twin} is the reference
against which the paper measures the price of second-order closure
(\cref{ch:attenuation-theorems}).

\section{The posterior mean as a denoiser}
\label{sec:ou:denoiser}

\begin{proposition}[The posterior mean is the best denoiser]
\label{prop:ou:mmse}
For every measurable $f\colon\Reals^{\nsites}\to\Reals^{\nsites}$ with $\E|f(\vx_t)|^2<\infty$,
\begin{equation}
  \E\bigl|\va-f(\vx_t)\bigr|^2=\E\bigl|\va-\pmean{\va}\bigr|^2+\E\bigl|\pmean{\va}-f(\vx_t)\bigr|^2 .
  \label{eq:ou:mmse}
\end{equation}
The first term on the right, written $\mathrm{MMSE}_t$, is the smallest mean squared
error any denoiser can reach; it is reached only by $f=\pmean{\va}$.
\end{proposition}

\begin{proof}
Expand $|\va-f|^2=|(\va-m)+(m-f)|^2$ with $m=\pmean{\va}$. The cross term is
$2\,\E\bigl\langle\va-m,\,m-f\bigr\rangle=2\,\E\bigl\langle\E[\va-m\mid\vx_t],\,m-f\bigr\rangle=0$,
because $m-f$ is a function of $\vx_t$ and $\E[\va-m\mid\vx_t]=0$.
\end{proof}

Two consequences organise the paper's measurements. First, any estimate
$\widehat{m}$ of the posterior mean converts into an estimate of the score and the
errors are proportional:
\begin{equation}
  \widehat{\score}=-\frac{\vx-e^{-t}\widehat{m}}{\Dt}
  \quad\Longrightarrow\quad
  \widehat{\score}-\score=\frac{e^{-t}}{\Dt}\bigl(\widehat{m}-\pmean{\va}\bigr),
  \label{eq:ou:score-mean-identity}
\end{equation}
an exact identity that \code{score\_mean\_identity\_residual} in
\code{src/markov\_score/metrics.py} checks to machine precision. Every estimator in the
paper, the tree-exact recursion, the fitted kernel, the LMMSE twin and the networks,
produces a posterior mean and is scored through \cref{eq:ou:score-mean-identity}.
Second, the energy of the score is tied to the MMSE.

\begin{proposition}[Score energy]
\label{prop:ou:energy}
$\E_{\Pt}|\score(\vx_t,t)|^2=\nsites/\Dt-(e^{-2t}/\Dt^2)\,\mathrm{MMSE}_t$.
\end{proposition}

\begin{proof}
By \cref{eq:ou:solution}, $-(\vx_t-e^{-t}\va)/\Dt=-\vz/\sqrt{\Dt}$, and by
\cref{thm:ou:tweedie} the score is its conditional expectation,
$\score(\vx_t,t)=\E[-\vz/\sqrt{\Dt}\mid\vx_t]$. Apply \cref{eq:ou:mmse} to the random
vector $-\vz/\sqrt{\Dt}$ in place of $\va$ and $f=\score$:
$\E|\vz|^2/\Dt=\E|\score|^2+\E|{-\vz/\sqrt{\Dt}}-\score|^2$. The last term is
$\E|\,e^{-t}(\va-\pmean{\va})/\Dt\,|^2=(e^{-2t}/\Dt^2)\mathrm{MMSE}_t$ by
\cref{eq:ou:score-mean-identity} applied to $\widehat m=\va$; and $\E|\vz|^2=\nsites$.
\end{proof}

At low noise the energy is close to $\nsites/\Dt$ and large; at high noise it tends to
$\nsites$, the energy of the score of $\gamma$. This is why the paper's pooled
estimand (\cref{ch:score-matching}) is dominated by its lowest levels.

\section{One noised view per sequence and the twelve-level schedule}
\label{sec:ou:protocol}

The paper observes $\nseq$ sequences $\va^{1},\dots,\va^{\nseq}$ and, in its own
estimator's interface (interface C of Section 5), sees each of them exactly once, at one
level $t_\mu$, as $\vx^{\mu}=e^{-t_\mu}\va^{\mu}+\sqrt{\Delta_{t_\mu}}\vz^{\mu}$; the clean
sequence, a second noised view, and the noise draw itself are never seen. The
sequences are split evenly across the twelve levels of \cref{tab:ou:levels}, so each
level receives $\nseq/12$ views (\code{\_one\_view\_groups} in
\code{experiments/e03\_baselines.py}; \code{one\_view=True} in
\code{experiments/protocol.py}). The likelihood of one view is then
$P_{t_\mu,\params}(\vx^{\mu})$, the density of \cref{eq:ou:convolution} under the
parameterised kernel, and the marginal log-likelihood of the paper's Section 4 is the
sum over views. The schedule is log-spaced so that the eleven ratios $t_{k+1}/t_k$ are
equal; on the SNR clock of \cref{eq:ou:snr} the levels run from an almost clean
observation (level 1) to almost pure noise (level 12), and the crossover scale
$\tstar$ of \cref{ch:attenuation-theorems} at the frozen correlation lies between the
ninth and the tenth level. The networks of the paper draw the level of each training
example uniformly from the same twelve values (\code{t\_levels} in
\code{src/markov\_score/nn\_torch.py}), so that their population objective is an
average over the schedule (\cref{ch:score-matching}).

\section{Worked example}
\label{sec:ou:example}

\begin{example}[A two-point prior at $t=\log2$]
\label{ex:ou:twopoint}
Let $\nsites=1$ and let $a=\pm1$ with probability $\tfrac12$ each, so that
$\Var(a)=1$. At $t=\log2$ we have $e^{-t}=\tfrac12$, $\Dt=1-\tfrac14=\tfrac34$ and
$\mathrm{SNR}=\tfrac13$. By \cref{eq:ou:convolution},
\[
  \Pt(x)=\tfrac12\gamma_{3/4}\bigl(x-\tfrac12\bigr)+\tfrac12\gamma_{3/4}\bigl(x+\tfrac12\bigr).
\]
The posterior of $a$ given $x$ puts mass proportional to $\exp[-(x\mp\tfrac12)^2/(2\cdot\tfrac34)]$
on $\pm1$; the ratio of the two weights is $\exp(2\cdot\tfrac12x/\tfrac34)=\exp(4x/3)$, so
\[
  \pmean{a}=\tanh\bigl(\tfrac23x\bigr),\qquad
  \score(x,t)=-\frac{x-\tfrac12\tanh(\tfrac23x)}{3/4}=-\tfrac43\Bigl(x-\tfrac12\tanh\tfrac{2x}{3}\Bigr).
\]
At $x=1$: $e^{4/3}=3.7937$, so $\tanh(2/3)=(3.7937-1)/(3.7937+1)=0.5828$,
$\pmean{a}=0.5828$ and $\score(1,t)=-\tfrac43(1-0.2914)=-0.9448$. The covariance twin
of this prior is $\normal01$, whose noised law is again $\normal01$; its score at $x=1$
is $\Sgauss=-1$ and its posterior mean is $e^{-t}x=0.5$. The true denoiser pulls
harder towards $+1$ than the linear one, by $0.083$, and the score gap
$\score-\Sgauss=+0.0552$ is, by \cref{eq:ou:score-mean-identity}, exactly
$(e^{-t}/\Dt)(0.5828-0.5)=\tfrac23\cdot0.0828$. As a check of \cref{cor:ou:gaussian}
on the twin: $-(x-\tfrac12\cdot\tfrac12x)/\tfrac34=-x$, as it must be.
\end{example}

\begin{example}[Reading the schedule]
\label{ex:ou:schedule}
At level 8 of \cref{tab:ou:levels}, $t=0.6769$, a site with clean value $a_i=2$ is
observed as $x_i=0.5082\cdot2+0.8612\,z_i=1.016+0.861\,z_i$: the signal has half its
amplitude and the noise has standard deviation $\sqrt{0.7417}=0.861$. At level 12 the
same site is observed as $0.0996+0.9988\,z_i$; its amplitude is one twentieth and the
SNR is $1/400$. The information a single view carries about the site's fourth
cumulant is, by the paper's Theorem 3, proportional to $e^{-8t}$, which between these
two levels falls by the factor $e^{-8(3-0.6769)}=e^{-18.58}\approx8.5\times10^{-9}$.
\end{example}

\where{The channel is equation (2) of the paper (Section 2, ``Data and channel'') and
Tweedie's identity is the first sentence of ``The score is a posterior mean''; the
Gaussian corollary is the covariance twin of the same section and of Proposition 1.
The twelve-level schedule is the frozen configuration of Appendix C, and the one-view
interface is interface C of Section 5. Code: the channel and its likelihood table are
\code{src/\allowbreak markov\_score/\allowbreak noising.py}; the LMMSE mean and score and the mean-to-score
conversion are in \code{src/\allowbreak markov\_score/\allowbreak theory.py}; the score-to-mean readout is in
\code{src/\allowbreak markov\_score/\allowbreak reverse.py}; the schedule is \code{t\_grid} in
\code{experiments/\allowbreak frozen\_config.py} and the one-view flag is in
\code{experiments/\allowbreak protocol.py}. Experiments: every registered experiment E01--E21
noises its data through the two lines of \cref{eq:ou:solution}.}

\paragraph{References.}
G.~E.~Uhlenbeck and L.~S.~Ornstein, 1930, ``On the theory of the Brownian motion'',
Physical Review 36, 823--841.
B.~\O ksendal, 2003, \emph{Stochastic Differential Equations: An Introduction with
Applications}, sixth edition, Springer.
H.~Robbins, 1956, ``An empirical Bayes approach to statistics'', Proceedings of the
Third Berkeley Symposium on Mathematical Statistics and Probability, vol.~1, 157--163.
B.~Efron, 2011, ``Tweedie's formula and selection bias'', Journal of the American
Statistical Association 106, 1602--1614.
Y.~Song, J.~Sohl-Dickstein, D.~P.~Kingma, A.~Kumar, S.~Ermon and B.~Poole, 2021,
``Score-based generative modeling through stochastic differential equations'',
International Conference on Learning Representations.
J.~Ho, A.~Jain and P.~Abbeel, 2020, ``Denoising diffusion probabilistic models'',
Advances in Neural Information Processing Systems 33.
